% --- Archon physics formalization source begin ---
% archon:physics
% archon:covers PhyXMiniProblems/problem_phyx_mini_0066.lean
% archon:source-report reports/phyx_mini/problem_phyx_mini_0066.source.json

\chapter{Physics problem phyx\_mini\_0066}
\label{ch:PhyXMiniProblems_problem_phyx_mini_0066}

\paragraph{Problem source.}
\#\# Question

Find the focal length of the glass lens in Figure

\#\# Answer choices

- A: A: -24cm

- B: B: -56cm

- C: C: -40cm

- D: D: -68cm

\#\# Figure caption (auxiliary; use the image as primary evidence)

The image depicts a diagram involving a concave lens. The lens is centrally positioned with its principal axis indicated by a horizontal line passing through it.

Key features include:

- The lens is shaded blue and has a concave shape.
- Two black dots are placed on the principal axis, equidistant from the lens on both sides, which likely represent points such as object or image positions.
- The distance from each black dot to the center of the lens is labeled "40 cm" with arrows pointing towards the lens.
- The arrows indicate the measurement of these distances along the principal axis.

The diagram appears to illustrate the lens's focal points or object/image positioning in relation to the lens.

\#\# Dataset metadata

Geometrical Optics; Physical Model Grounding Reasoning, Spatial Relation Reasoning

\paragraph{Recorded answer/context.}
C: C: -40cm

\paragraph{Figure/image path.}
/root/proposal\_for\_physic/science-mango/phyx\_mini\_run/phyx\_data/test\_image/66.png

\paragraph{Formalization target.}
create a compiling Lean file with sorry bodies at `PhyXMiniProblems/problem\_phyx\_mini\_0066.lean`. The Lean declarations must preserve the physical quantities, dimensions or dimensional roles, figure labels, governing-law hypotheses, and final relation expressed by this problem.
Use Mathlib/Physlib names found through LeanExplore where available. If a physics API is missing, introduce faithful local abstractions rather than scalar placeholder aliases.

\begin{theorem}[Physics formalization target]
\label{thm:physics:phyx_mini_0066:target}
\lean{PhyXMiniProblems.ProblemPhyXMini0066.problem_phyx_mini_0066}
\uses{def:physics:phyx-mini-0066:phyxminiproblems-problemphyxmini0066-lengthincentimeters, def:physics:phyx-mini-0066:phyxminiproblems-problemphyxmini0066-concavelenssetup, def:physics:phyx-mini-0066:phyxminiproblems-problemphyxmini0066-matchesconcavelensfigure, def:physics:phyx-mini-0066:phyxminiproblems-problemphyxmini0066-hasordinaryglassopticalparameters, def:physics:phyx-mini-0066:phyxminiproblems-problemphyxmini0066-satisfiesbiconcavelensclassification, def:physics:phyx-mini-0066:phyxminiproblems-problemphyxmini0066-satisfiesprincipalfocuslaw, def:physics:phyx-mini-0066:phyxminiproblems-problemphyxmini0066-matchesanswerexactly}
The assigned autoformalize agent should translate this physics problem into Lean declarations in the covered file, with theorem and lemma proof bodies written as `by sorry`.
\end{theorem}
\begin{proof}
This is an autoformalization task, not a proof task. Produce statements that can later be proved by the physics prover without weakening the physical model.
\end{proof}

% --- Archon declaration topology begin ---
\section{Declaration topology}
\begin{definition}[Length Quantity]
  \label{def:physics:phyx-mini-0066:phyxminiproblems-problemphyxmini0066-lengthquantity}
  \lean{PhyXMiniProblems.ProblemPhyXMini0066.LengthQuantity}
  A signed physical length whose value transforms coherently with the unit choice
\end{definition}
\begin{proof}
  This is the declared model definition of Length Quantity.
\end{proof}
\begin{definition}[centimeter Unit Choices]
  \label{def:physics:phyx-mini-0066:phyxminiproblems-problemphyxmini0066-centimeterunitchoices}
  \lean{PhyXMiniProblems.ProblemPhyXMini0066.centimeterUnitChoices}
  SI units with centimeters selected for length readouts
\end{definition}
\begin{proof}
  This is the declared model definition of centimeter Unit Choices.
\end{proof}
\begin{definition}[length In Centimeters]
  \label{def:physics:phyx-mini-0066:phyxminiproblems-problemphyxmini0066-lengthincentimeters}
  \lean{PhyXMiniProblems.ProblemPhyXMini0066.lengthInCentimeters}
  \uses{def:physics:phyx-mini-0066:phyxminiproblems-problemphyxmini0066-lengthquantity, def:physics:phyx-mini-0066:phyxminiproblems-problemphyxmini0066-centimeterunitchoices}
  The signed scalar readout of a physical length in centimeters
\end{definition}
\begin{proof}
  This is the declared model definition of length In Centimeters.
\end{proof}
\begin{definition}[Optical Medium]
  \label{def:physics:phyx-mini-0066:phyxminiproblems-problemphyxmini0066-opticalmedium}
  \lean{PhyXMiniProblems.ProblemPhyXMini0066.OpticalMedium}
  The optical media needed to classify a glass lens surrounded by air
\end{definition}
\begin{proof}
  This is the declared model definition of Optical Medium.
\end{proof}
\begin{definition}[Lens Material]
  \label{def:physics:phyx-mini-0066:phyxminiproblems-problemphyxmini0066-lensmaterial}
  \lean{PhyXMiniProblems.ProblemPhyXMini0066.LensMaterial}
  The material from which the depicted thin lens is made
\end{definition}
\begin{proof}
  This is the declared model definition of Lens Material.
\end{proof}
\begin{definition}[Lens Profile]
  \label{def:physics:phyx-mini-0066:phyxminiproblems-problemphyxmini0066-lensprofile}
  \lean{PhyXMiniProblems.ProblemPhyXMini0066.LensProfile}
  Axial cross-section profiles relevant to the depicted lens
\end{definition}
\begin{proof}
  This is the declared model definition of Lens Profile.
\end{proof}
\begin{definition}[Thin Lens Kind]
  \label{def:physics:phyx-mini-0066:phyxminiproblems-problemphyxmini0066-thinlenskind}
  \lean{PhyXMiniProblems.ProblemPhyXMini0066.ThinLensKind}
  The two paraxial optical classes of a thin lens
\end{definition}
\begin{proof}
  This is the declared model definition of Thin Lens Kind.
\end{proof}
\begin{definition}[Axial Location]
  \label{def:physics:phyx-mini-0066:phyxminiproblems-problemphyxmini0066-axiallocation}
  \lean{PhyXMiniProblems.ProblemPhyXMini0066.AxialLocation}
  The three concrete axial locations visible in the figure
\end{definition}
\begin{proof}
  This is the declared model definition of Axial Location.
\end{proof}
\begin{definition}[Principal Focus Side]
  \label{def:physics:phyx-mini-0066:phyxminiproblems-problemphyxmini0066-principalfocusside}
  \lean{PhyXMiniProblems.ProblemPhyXMini0066.PrincipalFocusSide}
  The two principal focal-point roles of a thin lens
\end{definition}
\begin{proof}
  This is the declared model definition of Principal Focus Side.
\end{proof}
\begin{definition}[Concave Lens Setup]
  \label{def:physics:phyx-mini-0066:phyxminiproblems-problemphyxmini0066-concavelenssetup}
  \lean{PhyXMiniProblems.ProblemPhyXMini0066.ConcaveLensSetup}
  \uses{def:physics:phyx-mini-0066:phyxminiproblems-problemphyxmini0066-lengthquantity, def:physics:phyx-mini-0066:phyxminiproblems-problemphyxmini0066-opticalmedium, def:physics:phyx-mini-0066:phyxminiproblems-problemphyxmini0066-lensmaterial, def:physics:phyx-mini-0066:phyxminiproblems-problemphyxmini0066-lensprofile, def:physics:phyx-mini-0066:phyxminiproblems-problemphyxmini0066-thinlenskind, def:physics:phyx-mini-0066:phyxminiproblems-problemphyxmini0066-axiallocation, def:physics:phyx-mini-0066:phyxminiproblems-problemphyxmini0066-principalfocusside}
  The physical lens and the labelled one-dimensional geometry of its figure. principalFocusMarker records the optical interpretation of the two black dots without identifying their distance with the unknown focal length in the data structure itself
\end{definition}
\begin{proof}
  This is the declared model definition of Concave Lens Setup.
\end{proof}
\begin{definition}[axial Separation In Centimeters]
  \label{def:physics:phyx-mini-0066:phyxminiproblems-problemphyxmini0066-axialseparationincentimeters}
  \lean{PhyXMiniProblems.ProblemPhyXMini0066.axialSeparationInCentimeters}
  \uses{def:physics:phyx-mini-0066:phyxminiproblems-problemphyxmini0066-lengthincentimeters, def:physics:phyx-mini-0066:phyxminiproblems-problemphyxmini0066-axiallocation, def:physics:phyx-mini-0066:phyxminiproblems-problemphyxmini0066-concavelenssetup}
  The unsigned separation of two labelled axis locations, in centimeters
\end{definition}
\begin{proof}
  This is the declared model definition of axial Separation In Centimeters.
\end{proof}
\begin{definition}[Matches Concave Lens Figure]
  \label{def:physics:phyx-mini-0066:phyxminiproblems-problemphyxmini0066-matchesconcavelensfigure}
  \lean{PhyXMiniProblems.ProblemPhyXMini0066.MatchesConcaveLensFigure}
  \uses{def:physics:phyx-mini-0066:phyxminiproblems-problemphyxmini0066-lengthincentimeters, def:physics:phyx-mini-0066:phyxminiproblems-problemphyxmini0066-concavelenssetup, def:physics:phyx-mini-0066:phyxminiproblems-problemphyxmini0066-axialseparationincentimeters}
  The qualitative labels and quantitative arrows read from the primary figure. The two 40 cm relations are retained independently, even though either one is sufficient to determine the focal-length magnitude
\end{definition}
\begin{proof}
  This is the declared model definition of Matches Concave Lens Figure.
\end{proof}
\begin{definition}[Has Ordinary Glass Optical Parameters]
  \label{def:physics:phyx-mini-0066:phyxminiproblems-problemphyxmini0066-hasordinaryglassopticalparameters}
  \lean{PhyXMiniProblems.ProblemPhyXMini0066.HasOrdinaryGlassOpticalParameters}
  \uses{def:physics:phyx-mini-0066:phyxminiproblems-problemphyxmini0066-concavelenssetup}
  The physical refractive-index branch for ordinary glass surrounded by air
\end{definition}
\begin{proof}
  This is the declared model definition of Has Ordinary Glass Optical Parameters.
\end{proof}
\begin{definition}[Satisfies Biconcave Lens Classification]
  \label{def:physics:phyx-mini-0066:phyxminiproblems-problemphyxmini0066-satisfiesbiconcavelensclassification}
  \lean{PhyXMiniProblems.ProblemPhyXMini0066.SatisfiesBiconcaveLensClassification}
  \uses{def:physics:phyx-mini-0066:phyxminiproblems-problemphyxmini0066-concavelenssetup}
  The standard paraxial classification of a biconcave glass lens in a lower-index ambient medium. This law determines only that the lens is diverging; it does not specify its focal-length magnitude
\end{definition}
\begin{proof}
  This is the declared model definition of Satisfies Biconcave Lens Classification.
\end{proof}
\begin{definition}[Satisfies Principal Focus Law]
  \label{def:physics:phyx-mini-0066:phyxminiproblems-problemphyxmini0066-satisfiesprincipalfocuslaw}
  \lean{PhyXMiniProblems.ProblemPhyXMini0066.SatisfiesPrincipalFocusLaw}
  \uses{def:physics:phyx-mini-0066:phyxminiproblems-problemphyxmini0066-lengthincentimeters, def:physics:phyx-mini-0066:phyxminiproblems-problemphyxmini0066-concavelenssetup, def:physics:phyx-mini-0066:phyxminiproblems-problemphyxmini0066-axialseparationincentimeters}
  The governing relation between the two principal focal points and signed focal length. Each principal-focus separation is the magnitude |f|, while the Cartesian sign convention assigns a negative focal length to a diverging lens. No numerical focal length is assumed here
\end{definition}
\begin{proof}
  This is the declared model definition of Satisfies Principal Focus Law.
\end{proof}
\begin{definition}[Answer Choice]
  \label{def:physics:phyx-mini-0066:phyxminiproblems-problemphyxmini0066-answerchoice}
  \lean{PhyXMiniProblems.ProblemPhyXMini0066.AnswerChoice}
  Labels of the four multiple-choice answers printed with the problem
\end{definition}
\begin{proof}
  This is the declared model definition of Answer Choice.
\end{proof}
\begin{definition}[answer Focal Length In Centimeters]
  \label{def:physics:phyx-mini-0066:phyxminiproblems-problemphyxmini0066-answerfocallengthincentimeters}
  \lean{PhyXMiniProblems.ProblemPhyXMini0066.answerFocalLengthInCentimeters}
  \uses{def:physics:phyx-mini-0066:phyxminiproblems-problemphyxmini0066-answerchoice}
  The signed focal-length readout printed for each answer, in centimeters
\end{definition}
\begin{proof}
  This is the declared model definition of answer Focal Length In Centimeters.
\end{proof}
\begin{definition}[Matches Answer Exactly]
  \label{def:physics:phyx-mini-0066:phyxminiproblems-problemphyxmini0066-matchesanswerexactly}
  \lean{PhyXMiniProblems.ProblemPhyXMini0066.MatchesAnswerExactly}
  \uses{def:physics:phyx-mini-0066:phyxminiproblems-problemphyxmini0066-lengthquantity, def:physics:phyx-mini-0066:phyxminiproblems-problemphyxmini0066-lengthincentimeters, def:physics:phyx-mini-0066:phyxminiproblems-problemphyxmini0066-answerchoice, def:physics:phyx-mini-0066:phyxminiproblems-problemphyxmini0066-answerfocallengthincentimeters}
  Exact agreement with one of the displayed whole-centimeter answers
\end{definition}
\begin{proof}
  This is the declared model definition of Matches Answer Exactly.
\end{proof}
\begin{lemma}[focal Length Magnitude eq forty]
  \label{lem:physics:phyx-mini-0066:phyxminiproblems-problemphyxmini0066-focallengthmagnitude-eq-forty}
  \lean{PhyXMiniProblems.ProblemPhyXMini0066.focalLengthMagnitude_eq_forty}
  \uses{def:physics:phyx-mini-0066:phyxminiproblems-problemphyxmini0066-lengthincentimeters, def:physics:phyx-mini-0066:phyxminiproblems-problemphyxmini0066-concavelenssetup, def:physics:phyx-mini-0066:phyxminiproblems-problemphyxmini0066-matchesconcavelensfigure, def:physics:phyx-mini-0066:phyxminiproblems-problemphyxmini0066-satisfiesprincipalfocuslaw}
  Either focal-point arrow in the figure, together with the principal-focus law, determines the magnitude of the focal length as 40 cm
\end{lemma}
\begin{proof}
  The conclusion follows from the governing relations and figure readouts named in the statement of focal Length Magnitude eq forty.
\end{proof}
% --- Archon declaration topology end ---
% --- Archon physics formalization source end ---
